\chapter{Learning as Optimization}

\section{Optimization Formulations of Learning}

In previous chapters, we formulated learning as the problem of minimizing empirical risk. We now examine this problem through the lens of optimization.

\medskip

Recall that a model is parameterized by $\theta \in \Theta$, and learning consists of solving
\[
\min_{\theta \in \Theta} \; \hat{R}_n(\theta)
\quad \text{where} \quad
\hat{R}_n(\theta) = \frac{1}{n} \sum_{i=1}^n \ell(f_\theta(x_i), y_i).
\]

\medskip

\noindent
\textbf{Key observation.}  
Learning reduces to minimizing a function over a high-dimensional parameter space.

\medskip

This formulation unifies a wide range of models:
\begin{itemize}
    \item Linear regression $\rightarrow$ convex quadratic optimization
    \item Logistic regression $\rightarrow$ convex optimization
    \item Neural networks $\rightarrow$ highly nonconvex optimization
\end{itemize}

\medskip

\noindent
\textbf{Structure of the objective.}
\begin{itemize}
    \item Sum of losses over data points
    \item Often smooth (differentiable)
    \item High-dimensional and possibly nonconvex
\end{itemize}

\medskip

\noindent
\textbf{Engineering constraint.}  
Exact minimization is rarely possible. Instead, we aim to find parameters that are \emph{good enough} within limited computational time.

\medskip

Thus, optimization is not merely a mathematical tool—it determines what models are practically usable.

\section{Gradient-Based Methods}

Most modern learning algorithms rely on gradient-based optimization.

\medskip

\noindent
\textbf{Gradient descent.}  
Given a differentiable objective $L(\theta)$, the update rule is
\[
\theta_{t+1} = \theta_t - \eta \nabla L(\theta_t),
\]
where $\eta > 0$ is the learning rate.

\medskip

\noindent
\textbf{Interpretation.}
\begin{itemize}
    \item The gradient $\nabla L(\theta)$ points in the direction of steepest increase
    \item Moving in the negative gradient direction decreases the loss
\end{itemize}

\medskip

\noindent
\textbf{Geometric view.}  
Gradient descent follows the local geometry of the loss surface, moving toward regions of lower value.

\medskip

\noindent
\textbf{Choice of learning rate.}
\begin{itemize}
    \item Too small $\rightarrow$ slow convergence
    \item Too large $\rightarrow$ instability or divergence
\end{itemize}

\medskip

\noindent
\textbf{Example (quadratic loss).}  
For a convex quadratic function, gradient descent converges to the global minimum under suitable conditions.

\medskip

\noindent
\textbf{Limitation.}  
Computing the full gradient requires summing over all data points:
\[
\nabla L(\theta) = \frac{1}{n} \sum_{i=1}^n \nabla \ell(f_\theta(x_i), y_i).
\]

For large datasets, this becomes computationally expensive.

\section{Stochastic Optimization}

To address the computational cost of full gradients, we use stochastic approximations.

\medskip

\noindent
\textbf{Stochastic Gradient Descent (SGD).}  
At each step, sample a data point (or mini-batch) and compute an approximate gradient:
\[
\theta_{t+1} = \theta_t - \eta \nabla \ell(f_\theta(x_{i_t}), y_{i_t}).
\]

\medskip

\noindent
\textbf{Mini-batch SGD.}
\[
\theta_{t+1} = \theta_t - \eta \frac{1}{|B_t|} \sum_{i \in B_t} \nabla \ell(f_\theta(x_i), y_i).
\]

\medskip

\noindent
\textbf{Key properties.}
\begin{itemize}
    \item Lower computational cost per iteration
    \item Introduces noise into the optimization process
\end{itemize}

\medskip

\noindent
\textbf{Insight.}  
The noise in SGD is not purely detrimental—it can help escape poor local minima and improve generalization.

\medskip

\noindent
\textbf{Tradeoff.}
\begin{itemize}
    \item Small batch size $\rightarrow$ noisy but cheap updates
    \item Large batch size $\rightarrow$ stable but expensive updates
\end{itemize}

\medskip

\noindent
\textbf{Practical methods.}
\begin{itemize}
    \item Momentum: accelerates convergence by smoothing updates
    \item Adaptive methods (e.g., Adam): adjust learning rates per parameter
\end{itemize}

\medskip

Stochastic optimization is the backbone of modern deep learning.

\section{Geometry of Loss Landscapes}

To understand optimization behavior, it is useful to consider the geometry of the loss function.

\medskip

\noindent
\textbf{Loss landscape.}  
The function
\[
L(\theta)
\]
defines a surface over parameter space.

\medskip

\noindent
\textbf{Critical points.}
\begin{itemize}
    \item Local minima: $\nabla L(\theta) = 0$, $L$ is locally minimal
    \item Saddle points: $\nabla L(\theta) = 0$, but not minima
\end{itemize}

\medskip

\noindent
\textbf{Convex vs nonconvex.}
\begin{itemize}
    \item Convex problems: all local minima are global
    \item Nonconvex problems: multiple local minima and saddle points
\end{itemize}

\medskip

Neural networks give rise to highly nonconvex loss landscapes, yet gradient-based methods often succeed in practice.

\medskip

\noindent
\textbf{Flat vs sharp minima.}
\begin{itemize}
    \item Flat minima: loss changes slowly around the solution
    \item Sharp minima: loss changes rapidly
\end{itemize}

\medskip

\noindent
\textbf{Observation.}  
Solutions found by SGD tend to lie in flatter regions, which are often associated with better generalization.

\medskip

\noindent
\textbf{High-dimensional effects.}  
In high dimensions:
\begin{itemize}
    \item Saddle points are more common than poor local minima
    \item Many directions may have near-zero curvature
\end{itemize}

\medskip

\noindent
\textbf{Geometric intuition.}  
Optimization can be viewed as navigating a complex landscape shaped by:
\begin{itemize}
    \item model architecture,
    \item data distribution,
    \item parameterization.
\end{itemize}

\section{Optimization and Generalization}

Optimization does more than minimize training loss—it influences generalization.

\medskip

\noindent
\textbf{Key idea.}  
Different optimization algorithms may converge to different solutions, even if they achieve similar training error.

\medskip

\noindent
\textbf{Implicit bias.}  
Gradient-based methods exhibit a preference for certain solutions among many possible minimizers.

\medskip

\noindent
\textbf{Example.}
\begin{itemize}
    \item In overparameterized linear models, gradient descent often finds minimum-norm solutions
\end{itemize}

\medskip

\noindent
\textbf{Implication.}  
Optimization is part of the modeling process—it implicitly regularizes the learned function.

\section{A Unifying Perspective}

We can now refine the learning process:

\begin{enumerate}
    \item Define a loss function based on data and model
    \item Formulate an optimization problem
    \item Apply gradient-based methods (often stochastic)
    \item Obtain a solution shaped by both the objective and the optimization dynamics
\end{enumerate}

\medskip

\noindent
\textbf{Three interacting aspects:}
\begin{itemize}
    \item \textbf{Objective} (what is minimized)
    \item \textbf{Algorithm} (how it is minimized)
    \item \textbf{Geometry} (structure of the loss landscape)
\end{itemize}

\medskip

Modern machine learning succeeds not only because good objectives are defined, but because efficient optimization methods can exploit their structure.

\section{Outlook}

This chapter introduced optimization as a central component of learning:
\begin{itemize}
    \item learning as minimizing empirical risk,
    \item gradient-based methods as the primary tool,
    \item stochasticity as both a necessity and an advantage.
\end{itemize}

\medskip

In the next chapter, we study generalization more carefully, addressing the fundamental question of why models that fit data well can still perform well on unseen examples.

\medskip

\noindent
\textbf{Guiding principle.}  
Learning is not only about what we optimize, but also about how we optimize it.